\ifdefined\gkver\else\def\gkver{student}\fi
\documentclass[\gkver]{gaokaozhenti}
\begin{document}

\chapter{2011年浙江理综}
%% paperType: 理综物理部分

\begin{enumerate}
\item
%% number: 14
%% paperName: 2011年浙江理综第 14 题 6 分
%% typeId: 1
%% type: 单选题
%% chapter: 运动和力
%% point: 牛顿第三定律
%% method: 无
%% score: 6
%% degree: 450
%% duplicateId: 0
%% body:
如图所示，甲、乙两人在冰面上“拔河”。两人中间位置处有一分界线，约定先使对方过分界线者为赢。若绳子质量不计，冰面可看成光滑，则下列说法正确的是\xzanswer{C}

\onepicture[w=5.5cm]{figs/14.png}

\fourchoices[answer=C]
{甲对绳的拉力与绳对甲的拉力是一对平衡力}
{甲对绳的拉力与乙对绳的拉力是作用力与反作用力}
{若甲的质量比乙大，则甲能赢得“拔河”比赛的胜利}
{若乙收绳的速度比甲快，则乙能赢得“拔河”比赛的胜利}

%% answer: C
%% memo:
\memoanswer{
甲对绳的拉力与绳对甲的拉力是作用力和反作用力，故 A 错误；甲对绳的拉力与乙对绳的拉力是一对平衡力，故 B 错误；若甲的质量比乙大，则甲的加速度比乙的小，可知乙先到分界线，故甲能赢得“拔河”比赛的胜利，故 C 正确；收绳速度的快慢并不能决定“拔河”比赛的输赢，故 D 错误。
}

\item
%% number: 15
%% paperName: 2011年浙江理综第 15 题 6 分
%% typeId: 1
%% type: 单选题
%% chapter: 原子和原子核
%% point: 天然放射性
%% method: 无
%% score: 6
%% degree: 650
%% duplicateId: 0
%% body:
关于天然放射现象，下列说法正确的是\xzanswer{D}

\fourchoices[answer=D]
{$\alpha$射线是由氦原子核衰变产生}
{$\beta$射线是由原子核外电子电离产生}
{$\gamma$射线是由原子核外的内层电子跃迁产生}
{通过化学反应不能改变物质的放射性}

%% answer: D
%% memo:
\memoanswer{
放射线是从原子核中释放出来的，通过化学反应并不能改变物质的放射性，故正确答案为 D。
}

\item
%% number: 16
%% paperName: 2011年浙江理综第 16 题 6 分
%% typeId: 1
%% type: 单选题
%% chapter: 交流电
%% point: 变压器
%% method: 无
%% score: 6
%% degree: 450
%% duplicateId: 0
%% body:
如图所示，在铁芯上、下分别绕有匝数 $n_{1}=800$ 和 $n_{2}=200$ 的两个线圈，上线圈两端与 $u=51\sin 314t\UV$ 的交流电源相连，将下线圈两端接交流电压表，则交流电压表的读数可能是\xzanswer{A}

\onepicture[h=2.8cm]{figs/16.png}

\fourchoices[answer=A]
{$2.0\UV$}
{$9.0\UV$}
{$12.7\UV$}
{$144.0\UV$}

%% answer: A
%% memo:
\memoanswer{
交流电源电压的有效值 $U_{1}=\frac{51}{\sqrt{2}}\UV$，如果看成理想的来处理有 $U_{1}/U_{2}=n_{1}/n_{2}$，$n_{1}=800$，$n_{2}=200$，解得 $U_{2}=9.0\UV$，故交流电压表的读数应该小于 $9.0\UV$，所以答案为 A。
}

\item
%% number: 17
%% paperName: 2011年浙江理综第 17 题 6 分
%% typeId: 1
%% type: 单选题
%% chapter: 光学
%% point: 光的折射
%% method: 无
%% score: 6
%% degree: 450
%% duplicateId: 0
%% body:
“B 超”可用于探测人体内脏的病变状况。下图是超声波从肝脏表面入射，经折射与反射，最后从肝脏表面射出的示意图。超声波在进入肝脏发生折射时遵循的规律与光的折射规律类似，可表述为 $\sin\theta_{1}/\sin\theta_{2}=v_{1}/v_{2}$（式中 $\theta_{1}$ 是入射角，$\theta_{2}$ 是折射角，$v_{1}$、$v_{2}$ 分别是超声波在肝外和肝内的传播速度），超声波在肿瘤表面发生反射时遵循的规律与光的反射规律相同，已知 $v_{2}=0.9v_{1}$，入射点与出射点之间的距离是 $d$，入射角为 $i$，肿瘤的反射面恰好与肝脏表面平行，则肿瘤离肝脏表面的深度 $h$ 为\xzanswer{D}

\onepicture[w=4.5cm]{\includesvg{figs/17}}

\fourchoices[answer=D]
{$\frac{9d\sin i}{\sqrt{100-81\sin^{2}i}}$}
{$\frac{d\sqrt{81-100\sin^{2}i}}{100\sin i}$}
{$\frac{d\sqrt{81-100\sin^{2}i}}{20\sin i}$}
{$\frac{d\sqrt{100-81\sin^{2}i}}{18\sin i}$}

%% answer: D
%% memo:
\memoanswer{
已知入射角为 $i$，设折射角为 $r$，根据题意有 $\tan r=\frac{d/2}{h}$、$\frac{\sin i}{\sin r}=\frac{v_{1}}{v_{2}}$，而 $v_{2}=0.9v_{1}$，解得 $h=\frac{d\sqrt{100-81\sin^{2}i}}{18\sin i}$。
}

\item
%% number: 18
%% paperName: 2011年浙江理综第 18 题 6 分
%% typeId: 2
%% type: 多选题
%% chapter: 光学
%% point: 光的干涉
%% method: 无
%% score: 6
%% degree: 650
%% duplicateId: 0
%% body:
关于波动，下列说法正确的是\xzanswer{BD}

\fourchoices[answer=BD]
{各种波均会发生偏振现象}
{用白光做单缝衍射与双缝干涉实验，均可看到彩色条纹}
{声波传播过程中，介质中质点的运动速度等于声波的传播速度}
{已知地震波的纵波速度大于横波速度，此性质可用于横波的预警}

%% answer: BD
%% memo:
\memoanswer{
只有横波才能发生偏振现象，故 A 错误；用白光做单缝衍射与双缝干涉，都可以观察到彩色条纹，故 B 正确；声波在传播过程中，介质中质点的速度并不等于声波的传播速度，故 C 错误；已知地震波的纵波波速大于横波波速，此性质可用于横波的预警，故 D 正确。
}

\item
%% number: 19
%% paperName: 2011年浙江理综第 19 题 6 分
%% typeId: 2
%% type: 多选题
%% chapter: 曲线运动
%% point: 人造地球卫星
%% method: 无
%% score: 6
%% degree: 450
%% duplicateId: 0
%% body:
为了探测 X 星球，载着登陆舱的探测飞船在该星球中心为圆心，半径为 $r_{1}$ 的圆轨道上运动，周期为 $T_{1}$，总质量为 $m_{1}$。随后登陆舱脱离飞船，变轨到离星球更近的半径为 $r_{2}$ 的圆轨道上运动，此时登陆舱的质量为 $m_{2}$，则\xzanswer{AD}

\fourchoices[answer=AD]
{X 星球的质量为 $M=\frac{4\pi^{2}r_{1}^{3}}{GT_{1}^{2}}$}
{X 星球表面的重力加速度为 $g_{x}=\frac{4\pi^{2}r_{1}}{T_{1}^{2}}$}
{登陆舱在 $r_{1}$ 与 $r_{2}$ 轨道上运动时的速度大小之比为 $\frac{v_{1}}{v_{2}}=\sqrt{\frac{m_{1}r_{2}}{m_{2}r_{1}}}$}
{登陆舱在半径为 $r_{2}$ 轨道上做圆周运动的周期为 $T_{2}=T_{1}\sqrt{\frac{r_{2}^{3}}{r_{1}^{3}}}$}

%% answer: AD
%% memo:
\memoanswer{
根据 $G\frac{Mm_{1}}{r_{1}^{2}}=m\frac{4\pi^{2}}{T_{1}^{2}}r_{1}$、$G\frac{Mm_{1}}{r_{2}^{2}}=m\frac{4\pi^{2}}{T_{2}^{2}}r_{2}$，可得 $M=\frac{4\pi^{2}r_{1}^{3}}{GT_{1}^{2}}$、$T_{2}=T_{1}\sqrt{\frac{r_{2}^{3}}{r_{1}^{3}}}$，故 A、D 正确；

登陆舱在半径为 $r_{1}$ 的圆轨道上运动的向心加速度 $a=r_{1}\omega_{1}^{2}=\frac{4\pi^{2}}{T_{1}^{2}}r_{1}$，此加速度与 X 星球表面的重力加速度并不相等，故 B 错误；根据 $G\frac{Mm}{r^{2}}=m\frac{v^{2}}{r}$，得 $v=\sqrt{\frac{GM}{r}}$，则 $\frac{v_{1}}{v_{2}}=\sqrt{\frac{r_{2}}{r_{1}}}$，故 C 错误。
}

\item
%% number: 20
%% paperName: 2011年浙江理综第 20 题 6 分
%% typeId: 2
%% type: 多选题
%% chapter: 磁场
%% point: 带电粒子在磁场中的运动
%% method: 无
%% score: 6
%% degree: 450
%% duplicateId: 0
%% body:
利用如图所示装置可以选择一定速度范围内的带电粒子。图中板 MN 上方是磁感应强度大小为 $B$、方向垂直纸面向里的匀强磁场，板上有两条宽度分别为 $2d$ 和 $d$ 的缝，两缝近端相距为 $L$。一群质量为 $m$、电荷量为 $q$，具有不同速度的粒子从宽度为 $2d$ 的缝垂直于板 MN 进入磁场，对于能够从宽度为 $d$ 的缝射出的粒子，下列说法正确的是\xzanswer{BC}

\onepicture[w=6cm]{\includesvg{figs/20}}

\fourchoices[answer=BC]
{粒子带正电}
{射出粒子的最大速度为 $\frac{qB(L+3d)}{2m}$}
{保持 $d$ 和 $L$ 不变，增大 $B$，射出粒子的最大速度与最小速度之差增大}
{保持 $d$ 和 $B$ 不变，增大 $L$，射出粒子的最大速度与最小速度之差增大}

%% answer: BC
%% memo:
\memoanswer{
由左手定则可判断粒子带负电，故 A 错误；由题意知：粒子的最大半径 $r_{\max}=\frac{L+3d}{2}$、粒子的最小半径 $r_{\min}=\frac{L}{2}$，根据 $r=\frac{mv}{qB}$，可得 $v_{\max}=\frac{qB(L+3d)}{2m}$、$v_{\min}=\frac{qBL}{2m}$，则 $v_{\max}-v_{\min}=\frac{3qBd}{2m}$，故可知 B、C 正确，D 错误。
}

\item
%% number: 21
%% paperName: 2011年浙江理综第 21 题 10 分
%% typeId: 4
%% type: 实验题
%% chapter: 运动和力
%% point: 验证牛顿第二定律
%% method: 无
%% score: 10
%% degree: 650
%% duplicateId: 0
%% body:
在“探究加速度与力、质量的关系”实验时，已提供了小车，一端附有定滑轮的长木板、纸带、带小盘的细线、刻度尺、天平、导线。为了完成实验，还须从下图中选取实验器材，其名称是 \tkanswer[5cm]{学生电源、电磁打点计时器、钩码、砝码（或电火花计时器、钩码、砝码）}（漏选或全选得零分）；并分别写出所选器材的作用 \tkanswer[5cm]{学生电源为电磁打点计时器提供交流电源；电磁打点计时器（电火花计时器）记录小车运动的位置和时间；钩码用以改变小车的质量；砝码用以改变小车受到的拉力的大小，还可以用于测量小车的质量。}

\onepicture[w=9cm]{figs/21.png}

%% answer: 
%% memo:
\memoanswer{
电磁打点计时器（电火花计时器）记录小车运动的位置和时间；钩码用以改变小车的质量；砝码用以改变小车受到的拉力的大小，还可以用于测量小车的质量。如果选电磁打点计时器，则需要学生电源，如果选电火花计时器，则不需要学生电源。
}

\item
%% number: 22
%% paperName: 2011年浙江理综第 22 题 10 分
%% typeId: 4
%% type: 实验题
%% chapter: 电路
%% point: 伏安法测电阻
%% method: 无
%% score: 10
%% degree: 450
%% duplicateId: 0
%% body:
在“探究导体电阻与其影响因素的定量关系”实验中，为了探究 3 根材料未知、横截面积均为 $S=0.20\Ummq$ 的金属丝 a、b、c 的电阻率，采用如图所示的实验电路，M 为金属丝 c 的左端点，O 为金属丝 a 的右端点，P 是金属丝上可移动的接触点。在实验过程中，电流表读数始终 $I=1.25\UA$。电压表读数 $U$ 随 OP 间距离 $x$ 的变化如下表：

\onepicture[w=6cm, align=right]{figs/22a.png}

\begin{center}\small
\begin{tblr}{|c|c|c|c|c|c|c|c|c|c|c|c|c|c|c|}
\hline
$x/\Umm$ & 600 & 700 & 800 & 900 & 1000 & 1200 & 1400 & 1600 & 1800 & 2000 & 2100 & 2200 & 2300 & 2400 \\
\hline
$U/\UV$ & 3.95 & 4.50 & 5.10 & 5.90 & 6.50 & 6.65 & 6.82 & 6.93 & 7.02 & 7.15 & 7.85 & 8.50 & 9.05 & 9.75 \\
\hline
\end{tblr}
\end{center}

\begin{enumerate}
	\item
	绘出电压表读数 $U$ 随 OP 间距离 $x$ 变化的图线；
	\item
	求出金属丝的电阻率 $\rho$，并进行比较。
\end{enumerate}

\onepicture[w=7cm, align=right]{figs/22b.png}

%% answer: 
\jdanswer{
\begin{enumerate}
	\item
	如图。

	\item
	$\rho_{a}=1.04\times10^{-6}\UOm$，$\rho_{b}=9.6\times10^{-7}\UOm$，$\rho_{c}=1.04\times10^{-6}\UOm$，金属丝 a 与 c 电阻率相同，远大于金属丝 b 的电阻率。

\end{enumerate}
}
%% memo:
\memoanswer{
（1）以 OP 间距离 $x$ 为横轴，以电压表读数 $U$ 为纵轴，描点、连线绘出电压表读数 $U$ 随 OP 间距离 $x$ 变化的图线。

\onepicture[w=7cm]{figs/22_an.png}

（2）根据电阻定律 $R=\rho\frac{l}{S}$ 可得 $\rho=R\frac{S}{l}=\frac{SU}{Il}$。

$\rho_{a}=\frac{(6.5-3.9)\times0.20\times10^{-6}}{1.25\times(1000-600)}\UOm=1.04\times10^{-6}\UOm$

$\rho_{b}=\frac{(7.1-6.5)\times0.20\times10^{-6}}{1.25\times(2000-1000)}\UOm=9.6\times10^{-7}\UOm$

$\rho_{c}=\frac{(9.7-7.1)\times0.20\times10^{-6}}{1.25\times(2400-2000)}\UOm=1.04\times10^{-6}\UOm$

通过计算可知，金属丝 a 与 c 电阻率相同，远大于金属丝 b 的电阻率。
}

\item
%% number: 23
%% paperName: 2011年浙江理综第 23 题 16 分
%% typeId: 6
%% type: 计算题
%% chapter: 电磁感应
%% point: 电磁感应能量分析
%% method: 无
%% score: 16
%% degree: 450
%% duplicateId: 0
%% body:
如图~\ref{2011浙江理综23a}所示，在水平面上固定有长为 $L=2\Um$、宽为 $d=1\Um$ 的金属“U”型导轨，在“U”型导轨右侧 $l=0.5\Um$ 范围内存在垂直纸面向里的匀强磁场，且磁感应强度随时间变化规律如图~\ref{2011浙江理综23b}所示。在 $t=0$ 时刻，质量为 $m=0.1\Ukg$ 的导体棒以 $v_{0}=1\Ums$ 的初速度从导轨的左端开始向右运动，导体棒与导轨之间的动摩擦因数为 $\mu=0.1$，导轨与导体棒单位长度的电阻均为 $\lambda=0.1\UOpm$，不计导体棒与导轨之间的接触电阻及地球磁场的影响（取 $g=10\Umsq$）。

\begin{enumerate}
	\item
	通过计算分析 4 s 内导体棒的运动情况；
	\item
	计算 4 s 内回路中电流的大小，并判断电流方向；
	\item
	计算 4 s 内回路产生的焦耳热。
\end{enumerate}

\twopicture[labelA=2011浙江理综23a, labelB=2011浙江理综23b, numstyle=chinese, align=right, h=3.0cm]
{figs/23a.png}
{figs/23b.png}

%% answer: 
\jdanswer{
\begin{enumerate}
	\item
	导体棒在 1 s 前做匀减速运动，在 1 s 后一直保持静止。

	\item
	0.2 A，电流方向是顺时针方向。

	\item
	0.04 J。

\end{enumerate}
}
%% memo:
\memoanswer{
（1）导体棒先在无磁场区域做匀减速运动，有 $-\mu mg=ma$，$v_{t}=v_{0}+at$，$x=v_{0}t+\frac{1}{2}at^{2}$，代入数据解得：$t=1\Us$，$x=0.5\Um$，导体棒没有进入磁场区域。

导体棒在 1 s 末已经停止运动，以后一直保持静止，离左端位置仍为 $x=0.5\Um$。

（2）前 2 s 磁通量不变，回路电动势和电流分别为 $E=0$，$I=0$，后 2 s 回路产生的电动势为 $E=\frac{\Delta\Phi}{\Delta t}=ld\frac{\Delta B}{\Delta t}=0.1\UV$，回路的总长度为 5 m，因此回路的总电阻为 $R=5\lambda=0.5\UO$，电流为 $I=\frac{E}{R}=0.2\UA$，根据楞次定律，在回路中的电流方向是顺时针方向。

（3）前 2 s 电流为零，后 2 s 有恒定电流，焦耳热为 $Q=I^{2}Rt=0.04\UJ$。
}

\item
%% number: 24
%% paperName: 2011年浙江理综第 24 题 20 分
%% typeId: 6
%% type: 计算题
%% chapter: 功和能
%% point: 交通工具启动
%% method: 无
%% score: 20
%% degree: 450
%% duplicateId: 0
%% body:
节能混合动力车是一种可以利用汽油及所储存电能作为动力来源的汽车。有一质量 $m=1000\Ukg$ 的混合动力轿车，在平直公路上以 $v_{1}=90\Ukmh$ 匀速行驶，发动机的输出功率为 $P=50\UW$。当驾驶员看到前方有 $80\Ukmh$ 的限速标志时，保持发动机功率不变，立即启动利用电磁阻尼带动的发电机工作给电池充电，使轿车做减速运动，运动 $L=72\Um$ 后，速度变为 $v_{2}=72\Ukmh$。此过程中发动机功率的 $\frac{1}{5}$ 用于轿车的牵引，$\frac{4}{5}$ 用于供给发电机工作，发动机输送给发电机的能量最后有 50\% 转化为电池的电能。假设轿车在上述运动过程中所受阻力保持不变。求：

\begin{enumerate}
	\item
	轿车以 $90\Ukmh$ 在平直公路上匀速行驶时，所受阻力 $f_{\text{阻}}$ 的大小；
	\item
	轿车从 $90\Ukmh$ 减速到 $72\Ukmh$ 过程中，获得的电能 $E_{\text{电}}$；
	\item
	轿车仅用其在上述减速过程中获得的电能 $E_{\text{电}}$ 维持 $72\Ukmh$ 匀速运动的距离 $L'$。
\end{enumerate}

%% answer: 
\jdanswer{
\begin{enumerate}
	\item
	$f_{\text{阻}}=2\times10^{3}\UN$

	\item
	$E_{\text{电}}=6.3\times10^{4}\UJ$

	\item
	$L'=31.5\Um$

\end{enumerate}
}
%% memo:
\memoanswer{
（1）汽车牵引力与输出功率的关系 $P=F_{\text{牵}}v$，将 $P=50\UW$，$v_{1}=90\Ukmh=25\Ums$，代入得

$F_{\text{牵}}=\frac{P}{v_{1}}=2\times10^{3}\UN$

当轿车匀速行驶时，牵引力与阻力大小相等，有

$F_{\text{阻}}=2\times10^{3}\UN$

（2）在减速过程中，注意到发动机只有 $\frac{1}{5}P$ 用于汽车的牵引，根据动能定理有 $\frac{1}{5}Pt-F_{\text{阻}}L=\frac{1}{2}mv_{2}^{2}-\frac{1}{2}mv_{1}^{2}$，代入数据得

$Pt=1.575\times10^{5}\UJ$

电源获得的电能为

$E_{\text{电}}=0.5\times\frac{4}{5}Pt=6.3\times10^{4}\UJ$

（3）根据题设，轿车在平直公路上匀速行驶时受到的阻力仍为 $F_{\text{阻}}=2\times10^{3}\UN$。此过程中，由能量转化及守恒定律可知，仅有电能用于克服阻力做功 $E_{\text{电}}=F_{\text{阻}}L'$，代入数据得 $L'=31.5\Um$。
}

\item
%% number: 25
%% paperName: 2011年浙江理综第 25 题 22 分
%% typeId: 6
%% type: 计算题
%% chapter: 电场
%% point: 带电粒子的直线运动
%% method: 无
%% score: 22
%% degree: 250
%% duplicateId: 0
%% body:
如图~\ref{2011浙江理综25a}所示，静电除尘装置中有一长为 $L$、宽为 $b$、高为 $d$ 的矩形通道，其前、后面板使用绝缘材料，上、下面板使用金属材料。图~\ref{2011浙江理综25b}是装置的截面图，上、下两板与电压恒定的高压直流电源相连。质量为 $m$、电荷量为 $-q$、分布均匀的尘埃以水平速度 $v_{0}$ 进入矩形通道，当带负电的尘埃碰到下板后其所带电荷被中和，同时被收集。通过调整两板间距 $d$ 可以改变收集效率 $\eta$。当 $d=d_{0}$ 时 $\eta$ 为 81\%（即离下板 $0.81d_{0}$ 范围内的尘埃能够被收集）。不计尘埃的重力及尘埃之间的相互作用。

\begin{enumerate}
	\item
	求收集效率为 100\% 时，两板间距的最大值为 $d_{m}$；
	\item
	求收集率 $\eta$ 与两板间距 $d$ 的函数关系；
	\item
	若单位体积内的尘埃数为 $n$，求稳定工作时单位时间下板收集的尘埃质量 $\frac{\Delta M}{\Delta t}$ 与两板间距 $d$ 的函数关系，并绘出图线。
\end{enumerate}

\twopicture[labelA=2011浙江理综25a, labelB=2011浙江理综25b, numstyle=chinese, align=right, h=3.6cm]
{figs/25a.png}
{figs/25b.png}

%% answer: 
\jdanswer{
\begin{enumerate}
	\item
	$d_{m}=0.9d_{0}$

	\item
	当 $d\leqslant0.9d_{0}$ 时，$\eta=100\%$；当 $d>0.9d_{0}$ 时，$\eta=0.81\left(\frac{d_{0}}{d}\right)^{2}$。

	\item
	当 $d\leqslant0.9d_{0}$ 时，$\frac{\Delta M}{\Delta t}=nmbdv_{0}$；当 $d>0.9d_{0}$ 时，$\frac{\Delta M}{\Delta t}=0.81nmbv_{0}\frac{d_{0}^{2}}{d}$，图线如图。

\end{enumerate}
}
%% memo:
\memoanswer{
（1）收集效率 $\eta$ 为 81\%，即离下板 $0.81d_{0}$ 的尘埃恰好到达下板的右端边缘，设高压电源的电压为 $U$，在水平方向有 $L=v_{0}t$，在竖直方向有 $0.81d_{0}=\frac{1}{2}at^{2}$，其中 $a=\frac{F}{m}=\frac{qE}{m}=\frac{qU}{md_{0}}$。当减小两板间距时，能够增大电场强度，提高装置对尘埃的收集效率。收集效率恰好为 100\% 时，两板间距为 $d_{m}$。如果进一步减小 $d$，收集效率仍为 100\%。因此，在水平方向有 $L=v_{0}t$，竖直方向有 $d_{m}=\frac{1}{2}a't^{2}$，其中 $a'=\frac{F'}{m}=\frac{qE'}{m}=\frac{qU}{md_{m}}$，联立可得 $d_{m}=0.9d_{0}$。

（2）通过前面的求解可知，当 $d\leqslant0.9d_{0}$ 时，收集效率 $\eta$ 为 100\%，当 $d>0.9d_{0}$ 时，设距下板 $x$ 处的尘埃恰好到达下板的右端边缘，此时有 $x=\frac{qUL^{2}}{2mdv_{0}^{2}}$，根据题意，收集效率为 $\eta=\frac{x}{d}$，联立可得 $\eta=0.81\left(\frac{d_{0}}{d}\right)^{2}$。

（3）稳定工作时单位时间下板收集的尘埃质量 $\frac{\Delta M}{\Delta t}=\eta nmbdv_{0}$，当 $d\leqslant0.9d_{0}$ 时，$\eta=1$，因此 $\frac{\Delta M}{\Delta t}=nmbdv_{0}$；当 $d>0.9d_{0}$ 时，$\eta=0.81\left(\frac{d_{0}}{d}\right)^{2}$，因此 $\frac{\Delta M}{\Delta t}=0.81nmbv_{0}\frac{d_{0}^{2}}{d}$，绘出的图线如下。

\onepicture[w=6cm]{figs/25_an.png}
}

\end{enumerate}

\end{document}
